%%% lorem.tex --- 
%% 
%% Filename: lorem.tex
%% Description: 
%% Author: Ola Leifler
%% Maintainer: 
%% Created: Wed Nov 10 09:59:23 2010 (CET)
%% Version: $Id$
%% Version: 
%% Last-Updated: Wed Nov 10 09:59:47 2010 (CET)
%%           By: Ola Leifler
%%     Update #: 2
%% URL: 
%% Keywords: 
%% Compatibility: 
%% 
%%%%%%%%%%%%%%%%%%%%%%%%%%%%%%%%%%%%%%%%%%%%%%%%%%%%%%%%%%%%%%%%%%%%%%
%% 
%%% Commentary: 
%% 
%% 
%% 
%%%%%%%%%%%%%%%%%%%%%%%%%%%%%%%%%%%%%%%%%%%%%%%%%%%%%%%%%%%%%%%%%%%%%%
%% 
%%% Change log:
%% 
%% 
%% RCS $Log$
%%%%%%%%%%%%%%%%%%%%%%%%%%%%%%%%%%%%%%%%%%%%%%%%%%%%%%%%%%%%%%%%%%%%%%
%% 
%%% Code:

\chapter{Discussion}
\label{cha:discussion}


\section{Results}
\label{sec:discussion-results}
In this section, we discuss the results from when applying the anomaly detection techniques on the different datasets and make comparisons between different algorithms.

\subsection{LogClustering}
The results clearly show that LogClustering performs the best on both datasets. This is likely because LogClustering uses a knowledge base with stored sequences from each cluster, which it can compare new sequences to. Looking at sequences works particularly well for system log files, since anomalies are often characterized by a typical event pattern. Additionally, since LogClustering is the only anomaly detection method which requires normal training data, it makes it easier to separate the normal clusters from anomalous in the evaluating phase. Although it might seem unfair to let LogClustering train on normal data while the other techniques don't, this report aims to find the best technique for online and offline scenarios. A offline real-world scenario would include knowing what normal logdata looks like, therefore it's not unreasonable to compare LogClustering to the others, giving each method what it needs to reach it's full potential. Further analyzing the results of LogClustering we see that it performs slightly better on the HDFS dataset than the BGL dataset. While this could be due to the different feature extraction techniques, it's not likely the main reason since the overall results show that offline techniques generate an F1-score of $80$\% on the HDFS data and $79$\% on the BGL data. These results are more likely due to the different dimensions, where $49$ unique event IDs were generated from the HDFS dataset, compared to BGL's $320$. More event IDs creates a high-dimensional and sparse event count matrix, making it difficult for the clustering to separate normal instances from anomalous, generating more false positives \cite{Shilin2017experiencereport}. This claim is further supported by the lower precision for when LogClustering is applied to the BGL data, compared to the $100$\% precision on the HDFS data.

\revrev{}{Lastly, when comparing LogClustering to the benchmark (SVM) on HDFS, we see that the F1-score is approximately 8 percentage points lower, but the precision performs better. However, this only indicates that LogClustering is precise, and that all data points classified as anomalies are indeed true anomalies. The lower recall of 81\% however indicates that it fails to capture all actual anomalies. The difference on BGL is larger, with a F1-score difference of approximately 16 percentage points, indicating that the higher complexity in BGL is harder to handle without labels.}
 
\subsection{PCA}
The PCA algorithm gave a poor precision result, only $59$\%, on the BGL dataset and made it the worst method among the three offline techniques. This may be because PCA identifies unusual data by measuring how much a point deviates from normal patterns and a single threshold might not be enough to capture a lot anomalies without many false positives \cite{Shilin2017experiencereport}. On the other hand, PCA produces a notably higher precision, $95$\%, for the HDFS dataset. The big difference between the algorithms performance on the datasets could be because the HDFS dataset was feature extracted using a session window method, whereas the BGL dataset used a sliding window approach. This introduces more variability and noise in the anomaly data, as further discussed in Section \ref{sec:discussion-method}. Our results also showed that it was easier to find anomalies in the BGL data (higher recall), this could also be due to the different feature extraction approaches. Overall the PCA algorithm provided higher precision on the HDFS dataset, resulting in a higher F1-score and therefore performing better on the HDFS data. \revrev{}{The PCA algorithm achieved an F1-score approximately 27 percentage points lower on the BGL data than the SVM benchmark technique. This significant difference indicates that PCA is far from perfect for this type of data. On the HDFS data, PCA was slightly closer to the SVM algorithm, now with a difference of approximately 15 percentage points in F1-score.}

\subsection{Isolation Forest}
Isolation Forest is the only offline algorithm that presents better results on the BGL data compared with the HDFS data. \revrev{}{Its performance on the HDFS dataset is significantly worse than the benchmark, with an \revrev{F1 score}{F1-score} that is 29 percentage points lower, compared to 19 percentage points lower on the BGL dataset.} The biggest difference seen when evaluating the results on the two datasets is that the recall is notably higher for the BGL dataset. Also, the precision is higher, but the difference is not as big. This indicates that even though BGL is feature extracted by a sliding window approach, the tree algorithm, Isolation Forest, manages to increase the precision slightly, which is a bit contradictory to what’s said in Section \ref{sec:discussion-method}. This leads to the conclusion that Isolation Forest handles noise that comes from the sliding window approach better than the other offline algorithms. It could also be that tree algorithms like Isolation Forest work better for sliding window approaches in an offline environment, but more research would need to be done on other tree algorithms to reach that conclusion.


\subsection{The online methods}
When comparing the online methods, the results lie fairly close to each other. Half-Space Trees (HST) performed best on the HDFS data, while One-Class SVM performed better on the BGL dataset. HDFS consists of fewer anomalous data points, and as mentioned in Section \ref{HST}, HST is particularly effective for datasets with few anomalous data points. This is also proven in our experimental results, where HST delivers an F1-score of $33$\% on HDFS compared to $21$\% on the BGL data. 

Furthermore, since BGL contains significantly more unique events than the HDFS dataset, this creates more unique combinations of events and anomalies. One-Class SVM can therefore achieve a better result on the BGL dataset than HST, since it has the ability to learn a compact decision boundary around the normal data in its feature space. \revrev{}{Also, it's kernel function helps dealing with complex and non-linear patterns, transforming input data into a higher dimensional space where linnear separation is possible.} This makes \revrev{a}{} One-Class SVM generalize better in scenarios with high variability, allowing it to detect anomalies that do not follow consistent patterns. In contrast, fewer unique events in the HDFS dataset makes HST more compatible, since the lower complexity results in a narrower range of event patterns and anomaly types. This suits HST and its way of isolating anomalies, which works well when anomalies clearly deviate from the normal data.

Although both have precision as its best performing metric in HDFS, and recall as its best performing metric on BGL, this is in line with what is suggested in Section \ref{sec:discussion-method}.

\subsection{Comparing online to offline anomaly detection}
The experimental results demonstrate that the offline approaches decisively outperform the online methods. These results comes with no surprises however, as the offline methods are training on the entire datasets, while the online methods are only trained on 80\% giving the offline methods more data to learn from. This also gives them the opportunity to score based on historical and future knowledge, while the online methods can only score based on historical knowledge. Additionally, the feature extraction will perform differently in the two scenarios. In the offline setting, complete feature vectors will be constructed throughout the entire dataset before any scoring takes place. Giving the anomaly detection model access to richer and more context-aware features. The online scenario on the other hand, must process logs incrementally, scoring row by row as the new data arrives. Therefore the feature vectors are updated dynamically, sometimes with incomplete contextual information. This fundamental limitation contributes to the performance gap between offline and online settings. 

\section{Method}
\label{sec:discussion-method}
Our feature extraction methods differed between the two datasets. The HDFS dataset included block IDs for each log entry and that allowed for a session window approach to be used. Since each log row was associated with a specific block ID, we could label entire blocks as either anomalous or normal. In contrast, the BGL dataset didn't have such identifiers, hindering the use of a session-based approach. Instead, we applied a sliding window method, where windows were classified as anomalous if they contained one or more anomalous log entries.

This difference in feature extraction had a considerable impact on performance. The session window approach led to higher precision, as anomalies were easier to identify\revrev{ .}{, likely because the vectors were grouped by \textit{block\_id}, which were already labeled as either anomalous or normal. This decreases the chances of false positives.} On the other hand, the sliding window method resulted in higher recall, since several anomalies and normals could be in the same window. This caused some windows to be labeled as anomalous even if they included events that were actually normal. \revrev{}{This leads to a higher recall since more actual anomalies are captured, but lower precision due to an increased number of false positives}. 

For the sliding window approach applied to the offline algorithms on the BGL dataset, we found that the step size had minimal impact on the results. The window size, on the other hand, had a significant effect on the results. This is likely because the step size only controls how frequently a new window is generated, without changing much of the actual content in each window. However, changes in window size affect what data is included within each window, which can be the determining factor in whether a window is classified as anomalous or normal.

Even though the largest window size (eight hours) consistently showed the best results for the offline algorithms, it also comes with some negative effects. A high window size means that we know that at least one anomaly exists within the window but we don't know exactly which log entry is anomalous or how many anomalies are present. As a result, manual inspection is required to find the anomalies, and a larger window leads to more log entries to examine. In summary, a larger window size improves the performance but increases the effort and time needed for analysis.

\subsection{Limitations}
Our results showed a notable performance gap between the online and offline algorithms. One key reason for this difference is that the online algorithms required longer runtime, which limited our ability to tune their parameters to the same extent as we did for the offline algorithms. For example, with Half-Space Trees the ideal would be to experiment with the number of trees and their height, but increasing these parameters made the algorithm unable to run within a reasonable time on our datasets and hardware. Likewise, One-Class SVM had a long runtime even with default settings, which made it difficult to work iteratively and explore parameter adjustments. Additionally, we were not able to adjust the window sizes for the online algorithms on the BGL dataset. The online techniques became very slow when modifying the window size, which prevented us from testing different window sizes. 

We provided the different algorithms with event IDs based on log templates. It means that the algorithms only had access to one type of parameter, the amount of times a specific event occurred. It could be beneficial to provide more parameters, such as size, timestamp and addresses, for better results and making it easier for the algorithm to detect anomalies. 

\section{The work in a wider context}
\label{sec:work-wider-context}
The collection and analysis of log data could raise some ethical questions due to different reasons. Firstly, using machine learning approaches for anomaly detection studies user behavior, and it should be made sure that the work follows laws and regulations such as GDPR, and the processing of data should ensure privacy for individuals \cite{Devineni2023datasecurityandintegrity}. Also, the suggested machine learning models almost always find some false positives. This can lead to falsely accusing a user for abnormal behavior, creating ethical concerns. Lastly, the data handled could be considered sensitive, and as Devineni et al.\cite{Devineni2023datasecurityandintegrity} claims, it is critical that data can be anonymized or receives secure encryption. However, fighting these challenges without affecting the results of the machine learning techniques is not self-evident.


%%%%%%%%%%%%%%%%%%%%%%%%%%%%%%%%%%%%%%%%%%%%%%%%%%%%%%%%%%%%%%%%%%%%%%
%%% lorem.tex ends here

%%% Local Variables: 
%%% mode: latex
%%% TeX-master: "demothesis"
%%% End: 
